%% hopfion-framework/chemistry/phi_spiral.tex
%% Sources: main_paper11.tex
%% phi-spiral periodic table, tower levels, Z_eff(O), Rydberg as tower level
%% Auto-generated from project files. Edit source papers to update.

%════════════════════════════════════════════════════════════════════
% phi-Spiral Definitions — Paper XI
%════════════════════════════════════════════════════════════════════
\begin{definition}[$\vp$-tower level]
\label{P11:def:tower_level}
For any energy scale $E > 0$, its \emph{$\vp$-tower level} is
\begin{equation}
  \nlvl{E} \;=\; -\frac{\ln(E/\Lcond)}{2\ln\vp}.
  \label{P11:eq:tower_level}
\end{equation}
Equivalently, $E = \Lcond\,\vp^{-2\,n(E)}$.
Level $n=0$ is the condensate ground state $\Lcond$;
increasing $|n|$ corresponds to higher energy.
\end{definition}
\status{published}
\source{P11:def:tower_level}

\begin{remark}[Quantum origin of the tower spacing]
\label{P11:rem:quantum_tower_origin}
The spacing $\vp^2$ between successive tower levels is not merely a
display choice: Paper~X~\cite{Paper10} (Theorem~6.3 therein) proves that
the quantised condensate Hamiltonian has spectrum $E_n = \vp^{2n}E_0$
and that the $q$-deformation parameter of the resulting Macfarlane--Biedenharn
oscillator algebra is exactly $q = \vp^2$.
The tower of Table~\ref{P11:tab:levels} is therefore the semiclassical limit
of a quantum spectrum, with $\vp^2$ being the quantum of the condensate ladder.
\end{remark}
\status{published}
\source{P11:rem:quantum_tower_origin}

\begin{definition}[$\vp$-spiral]
\label{P11:def:spiral}
The \emph{$\vp$-spiral} is the curve
$(\theta(n), r(n))$ in cylindrical coordinates defined by
\begin{equation}
  \theta(n) = \frac{2\pi\,|n|}{N_{\rm arms}},
  \qquad
  r(n) = R_0\,\log(1 + 0.8\,|n|),
  \qquad
  y(n) = |n|\cdot\delta y,
  \label{P11:eq:spiral}
\end{equation}
where $N_{\rm arms} = 6$ is the number of tower levels per revolution,
$R_0$ is a display scale factor, and $\delta y$ is a vertical rise per
level.
The logarithmic radial coordinate $r = R_0\log(1+0.8|n|)$ compresses
the 77-level range from $\Lcond$ to $\MPl$ into a finite viewing
region while preserving the relative angular positions that encode
the physics.
\end{definition}
\status{published}
\source{P11:def:spiral}

\begin{remark}[Choice of $N_{\rm arms}$]
\label{P11:rem:Narms_choice}
The value $N_{\rm arms} = 6$ is a display choice, not a physical one.
Any value gives the same physics; the angular groupings of elements
and particles change as $N_{\rm arms}$ varies, making the spiral an
interactive diagnostic tool.  The interactive Three.js visualisation
(Appendix~\ref{P11:app:code}) includes a live slider for $N_{\rm arms}$.
\end{remark}
\status{published}
\source{P11:rem:Narms_choice}

%════════════════════════════════════════════════════════════════════
% Tower Level Results — Paper XI
%════════════════════════════════════════════════════════════════════
\begin{corollary}[Rydberg as derived tower level]
\label{P11:cor:rydberg}
Within the density-feedback Hopfion framework,
the Rydberg energy $\ERy = \aem^2 m_e/2$ is a
\emph{calculable, parameter-free} tower level:
\begin{equation}
  \boxed{\nlvl{\ERy}
  \;=\; n(m_e) + \frac{\ln(2/\aem^2)}{2\ln\vp}
  \;=\; -12.102},
  \label{P11:eq:nRy}
\end{equation}
and the Rydberg energy itself is
\begin{equation}
  \boxed{
    \ERy
    \;=\; \frac{\aem^2\,m_e}{2}
    \;=\; 13.606\,\mathrm{eV}
    \;=\; \Lcond\,\vp^{24.204},
  }
  \label{P11:eq:Ry_value}
\end{equation}
where every factor is a framework output:
$\aem^{-1} = 360/\vp^2 - 3/(2\pi) + 1/(9\vp^6) - 1/(36\pi\vp^6)= 137.036$
(Paper~IV, Theorem~5.3 and Corollary~8.11, four-term, $5\times10^{-7}\%$)) and
$m_e = \Lcond\,\vp^{20}\,e^{(1600\pi-3)/400}\cdot(1+O(R_0^{-2}))$
(Paper~IV, Verlinde cancellation~\cite{Paper4}).
\end{corollary}
\status{published}
\source{P11:cor:rydberg}

\begin{remark}[Physical interpretation]
\label{P11:rem:rydberg_physical}
The Rydberg is the natural unit of atomic energy — the binding energy
of the hydrogen 1s electron in the Bohr model,
$\ERy = m_e \aem^2/2 = 1\,\mathrm{Ry}$.
Corollary~\ref{P11:cor:rydberg} says that this unit does not float freely:
it is pinned to the $\vp$-tower at level $-12.102$ by the same
condensate that fixes $\TCMB$ and $m_e$.
The Bohr radius $a_0 = 1/(\aem m_e)$ and all atomic length scales
follow, pinned to the same condensate.
Chemistry, in this sense, is \emph{downstream} of the condensate.

Since Paper~X~\cite{Paper10} establishes the full canonical quantisation
of the condensate (Hilbert space $\mathcal{H}=L^2(\mathcal{C}_Q^\mathrm{red},d\mu)$,
Born rule, and spectral gap), this statement holds at the quantum level:
$\aem$ and $m_e$ are outputs of a well-defined quantum field theory,
and the Rydberg is a calculable tower level of that quantum theory.
\end{remark}
\status{published}
\source{P11:rem:rydberg_physical}

\begin{proposition}[Red-light threshold as pure tower level]
\label{P11:prop:visible}
The red edge of the visible electromagnetic spectrum
($\approx 1.77\,\mathrm{eV}$) coincides with the
pure tower level at $|n|=20$:
\begin{equation}
  \Lcond\,\vp^{20}
  \;=\; 1.1895\times10^{-4}\,\mathrm{eV}\times\vp^{20}
  \;=\; 1.799\,\mathrm{eV},
  \label{P11:eq:red_light}
\end{equation}
accurate to $1.7\%$.
\end{proposition}
\status{published}
\source{P11:prop:visible}

\begin{remark}[Two uses of level 20]
\label{P11:rem:level20_dual_use}
The same level index $n=20$ appears in two different formulas:
\begin{align}
  \text{Red-light threshold:} &\quad \Lcond\,\vp^{20} = 1.799\,\mathrm{eV}, \\
  \text{Electron mass:} &\quad m_e = \Lcond\,\vp^{20}\,e^{4\pi-3/400}\cdot(1+O(R_0^{-2})),
\end{align}
where $e^{4\pi-3/400} \approx 2.87\times10^5$.
The pure tower level at $n=20$ is the red-light threshold.
The physical electron mass sits at the same tower index but is lifted
by the WZW Chern--Simons exponential $e^{4\pi}$ from this threshold
by a factor of $\approx 286\,000$.
In other words: the optical window ``knows about'' the electron mass
tower level, but not its WZW lift.
\end{remark}
\status{published}
\source{P11:rem:level20_dual_use}

\begin{remark}[What the spiral measures]
\label{P11:rem:spiral_measures}
Two elements at the same tower level have outermost electrons
bound with the same energy.  This is a statement about their
effective nuclear charge, independent of their position in
the conventional periodic table.  The spiral thus provides a
binding-energy organisation of chemistry, complementary to
the electron-configuration organisation of the standard table.
\end{remark}
\status{published}
\source{P11:rem:spiral_measures}

\begin{proposition}[Level-1: $R(\mathrm{O})\approx1$ from the eight-electron problem]
\label{P11:prop:R_equals_1}
The dimensionless ratio for the 2p ionisation of oxygen in the
$1s^22s^22p^4\,{}^3P$ ground state,
\begin{equation}
  R(\mathrm{O}) \;\equiv\; \frac{\IE(\mathrm{O})}{\ERy},
  \label{P11:eq:R_equals_1}
\end{equation}
is fixed by $Z=8$ and the configuration alone, and equals unity to the
few-per-cent accuracy of an analytic orbital treatment:
$\boxed{R(\mathrm{O}) = 1 + O(\text{few}\,\%)}$.  The measured value is
$R(\mathrm{O}) = 1.000\,90$; the residual $9.04\times10^{-4}$ excess is
$2p^4$ correlation energy --- itself condensate-downstream through the
framework-fixed Hamiltonian, but not a coupling shift beyond it
(Remark~\ref{P11:rem:oalpha_excess}).
\end{proposition}
\status{published}
\source{P11:prop:R_equals_1}

\begin{remark}[The $2p^4$ pairing penalty and the nitrogen--oxygen anomaly]
\label{P11:rem:pairing_penalty}
Slater's rules carry no exchange or spin dependence, so
equation~\eqref{P11:eq:IE_slater_total} cannot see the effect that
brings oxygen's first ionisation energy below nitrogen's despite the
larger nuclear charge: in $2p^4$ the fourth electron must doubly
occupy an orbital already singly occupied, at the cost of the pair
repulsion absent in the half-filled $2p^3$ shell.  That anomaly is
measured directly,
\begin{equation}
  \IE(\mathrm{N}) - \IE(\mathrm{O})
  = 14.534 - 13.618 = 0.916\,\mathrm{eV},
  \label{P11:eq:NO_anomaly}
\end{equation}
and the Condon--Shortley angular algebra~\cite{Condon1935} reproduces
it: with the angular coefficient $3/25$ of $K_{2p,2p}$ and the hydrogenic
radial integral $F^2(2p,2p) = (45/256)Z_\mathrm{eff}\ERy$
of~\eqref{P11:eq:slater_condon} evaluated at $Z_\mathrm{eff}=2$,
\begin{equation}
  \tfrac{3}{2}K_{2p,2p} = \tfrac{81}{2560}\,Z_\mathrm{eff}\,\ERy
  = 0.861\,\mathrm{eV},
  \label{P11:eq:pairing_K}
\end{equation}
within $6\%$ of~\eqref{P11:eq:NO_anomaly}.  The pairing penalty is
therefore the physical effect that lowers $\IE(\mathrm{O})$ onto the
Rydberg scale, and it is of the right size.

It does not combine additively with~\eqref{P11:eq:IE_slater_total} to
give a sharper prediction.  The fitted-screening estimate lies $0.55\,$eV above the measured
value and the pairing term is $0.86\,$eV, so subtracting one from the
other returns $13.31\,\mathrm{eV}$, now $2.3\%$ \emph{below}
experiment: the two parametrisations are not
disjoint contributions, since Slater's electrostatic constants are
fitted to configurations in which pairing is already present on
average.  What the two establish jointly is that ordinary atomic
structure brackets $\IE(\mathrm{O})$ around $\ERy$ from both sides at
the few-per-cent level --- the coincidence is a real feature of the
$1s^22s^22p^4$ configuration, not an accident of arithmetic --- while
neither route approaches the precision of the observed agreement.
\end{remark}
\status{published}
\source{P11:rem:pairing_penalty}

\begin{remark}[Precision of the H--O coincidence]
\label{P11:rem:sigma_xc_origin}
Equation~\eqref{P11:eq:IE_slater_total} accounts for
$R(\mathrm{O})=1$ to four per cent, $|\Delta n| = 0.041$.  The measured
coincidence is far tighter: $\IE(\mathrm{H})=13.598\,\mathrm{eV}$ and
$\IE(\mathrm{O})=13.618\,\mathrm{eV}$ differ by $0.15\%$,
$|\Delta n| = 0.0015$, some twenty-seven times closer than the
screening calculation predicts.  That excess precision is not derived
here, and is the substance of the near-degeneracy.

Two accounts of it are unavailable.  A decomposition
$\sigma_\mathrm{tot} = \sigma_\mathrm{Slater} + \sigma_\mathrm{xc}$
cannot supply it, because $\sigma_\mathrm{xc}$ would be defined as the
residual $\sigma_\mathrm{tot}-\sigma_\mathrm{Slater}$, making its
closure an identity rather than a check --- and, by the ceiling
argument above, no admissible $\sigma_\mathrm{xc}$ reaches
$\sigma_\mathrm{tot}=6.00$ in any case.  Nor does the $2p^4$ exchange
penalty account for the gap quantitatively: with the Condon--Shortley
angular coefficient $3/25$ and the hydrogenic
$F^2(2p,2p)=(45/256)Z_\mathrm{eff}\ERy$ at $Z_\mathrm{eff}=2$, the
exchange term is $(3/2)K_{2p,2p} = 0.86\,\mathrm{eV}$, which at
$\mathrm{d}\IE/\mathrm{d}\sigma = 2Z_\mathrm{eff}\ERy/n_\mathrm{qn}^2
= 13.6\,\mathrm{eV}$ per unit of screening corresponds to
$\Delta\sigma \approx 0.08$.

What the framework contributes is unaffected: the Rydberg energy sits
at tower level $-12.102$ with no free parameters once $\aem$ and $m_e$
are fixed (Corollary~\ref{P11:cor:rydberg}), and hydrogen sits there by
definition.  That oxygen sits within $0.0015$ of the same level is an
atomic-structure fact, approximately explained by ordinary screening
and exactly explained by neither.
\end{remark}
\status{published}
\source{P11:rem:sigma_xc_origin}

\begin{remark}[Why the $O(\aem)$ excess is not a condensate correction]
\label{P11:rem:oalpha_excess}
The measured first ionisation energy of oxygen exceeds the Rydberg
energy by
\begin{equation}
  \frac{\IE(\mathrm{O})}{\ERy} - 1 = 9.04\times10^{-4},
  \label{P11:eq:IE_O_excess}
\end{equation}
an $O(\aem)$ quantity, and the condensate value $\aem = 1/137.036$
makes
\begin{equation}
  \frac{k\,\aem}{8\pi}\bigg|_{k=3} = \frac{3\aem}{8\pi} = 8.71\times10^{-4}
  \label{P11:eq:IE_O_closed}
\end{equation}
a candidate closed form for it, reproducing the
excess~\eqref{P11:eq:IE_O_excess} to $3.7\%$ (a residual of $0.003\%$
of the total ionisation energy).  The coefficient has a motivating
reading: $k/(8\pi) = k/(4\pi\,\Delta Q)$ with $\Delta Q = 2$, the
Chern--Simons barrier of a two-unit topological transition being
$4\pi\,\Delta Q$ (Paper~XVI~\cite{Paper16}, where $\Delta Q = 3-1 = 2$
gives the $8\pi$ spoke of the constituent quark mass), and removing one
electron from the lepton-sector $\QH=2$ ground state
(Paper~I~\cite{Paper1}) is a $\Delta Q = 2$ event --- mirroring the
Chern--Simons factor $1/(4\pi)$ of the electron-mass formula of
Paper~IV~\cite{Paper4}.

The reading is nonetheless excluded as a physical mechanism by its own
universality.  Every single-electron ionisation is a $\Delta Q = 2$
event, so the correction would apply to hydrogen identically; but
$\IE(\mathrm{H})$ is accounted for by the reduced-mass factor alone,
$\ERy\,(1+m_e/m_p)^{-1} = 13.5983\,\mathrm{eV}$ against the measured
$13.5984\,\mathrm{eV}$ (residual $11\,\mathrm{ppm}$), and adding
$3\aem/(8\pi)\,\ERy = 11.9\,\mathrm{meV}$ would spoil this by
$860\,\mathrm{ppm}$, eighty times the residual.  A genuine QED
radiative shift at $Z=8$ is of order the Lamb scale
($\sim10^{-6}\,\mathrm{eV}$), three orders below the $11\,\mathrm{meV}$
here.  The excess~\eqref{P11:eq:IE_O_excess} is therefore many-body
structure --- the $2p^4$ correlation energy that separates Koopmans
from $\Delta$SCF (Proposition~\ref{P11:prop:R_equals_1}).  It is
condensate physics in the same parametric sense as the entire Coulomb
Hamiltonian, which is fixed once $\aem$ and $m_e$ are
(Step~1); what the hydrogen test excludes is the stronger claim that
the condensate adds a coupling shift \emph{beyond} that Hamiltonian,
universal across elements.  The agreement with $3\aem/(8\pi)$, closer
than the alternative $2\aem/(5\pi)$ ($2.8\%$) only by chance, is a
coincidence of the right order, not such a shift.  Whether the
condensate's structure instead surfaces \emph{through} the many-body
correlation dynamics --- which vanish for hydrogen and so evade this
test --- is a separate and open question, of the same kind as the
cooperative-binding bound of Section~\ref{P11:sec:open}.
\end{remark}
\status{draft}
\source{P11:rem:oalpha_excess}

\begin{remark}[What is fixed and what is not]
\label{P11:rem:two_level}
The oxygen placement separates into two statements of very different
standing:
\begin{enumerate}[noitemsep]
\item \emph{Level~1 (atomic structure, $O(\aem^0)$):}
  the ratio $\IE(\mathrm{O})/\ERy$ is fixed by $Z=8$ and the
  $1s^22s^22p^4$ configuration through the exact Slater--Condon
  angular algebra, but only to the tens-of-per-cent accuracy of an
  orbital treatment (Koopmans $+26\%$, $\Delta$SCF $-13\%$; the gap is
  correlation energy).  It gives $\IE(\mathrm{O}) \approx \ERy$, not
  $Z_\mathrm{eff}=2$ exactly.
\item \emph{Level~2 ($O(\aem^1)$):}
  the residual excess $9.04\times10^{-4}$ is coincidentally near
  $3\aem/(8\pi)$ but is not a condensate correction: the same form
  applied to hydrogen fails by eighty times its own residual
  (Remark~\ref{P11:rem:oalpha_excess}), so the excess is
  correlation energy, not a WZW radiative shift.
\end{enumerate}
What the framework contributes is Level~0: the Rydberg energy itself
sits at tower level $-12.102$ with no free parameters
(Corollary~\ref{P11:cor:rydberg}), and hydrogen sits there by
definition.  Oxygen's proximity to that level is atomic structure,
approximately explained and precisely unexplained:
\[
  \TCMB,\, m_e
  \;\xrightarrow{\text{Papers I--X}}\;
  \aem,\, m_e
  \;\xrightarrow{\text{Cor.~\ref{P11:cor:rydberg}}}\;
  n(\ERy) = -12.102
  \;\xrightarrow{\text{atomic structure}}\;
  \IE(\mathrm{O}) \approx \ERy .
\]
\end{remark}
\status{published}
\source{P11:rem:two_level}

\begin{remark}[Framework origin of the coefficient $k/(8\pi)$]
\label{P11:rem:k8pi_origin}
The candidate coefficient $k/(8\pi)$ reads as
$k/(4\pi\,\Delta Q)$ with $\Delta Q = 2$:
\begin{itemize}[noitemsep]
\item The denominator $4\pi\,\Delta Q = 8\pi$ is the Chern--Simons
  barrier of a two-unit topological transition
  (Paper~XVI~\cite{Paper16}, where $\Delta Q = 3-1 = 2$ gives the
  $8\pi$ spoke of the constituent quark mass), and single-electron
  ionisation of the $\QH=2$ lepton ground state is a $\Delta Q = 2$
  event.
\item The numerator $k=3$ is the WZW level of the $\mathrm{SU}(2)_3$
  condensate (Paper~I).
\end{itemize}
This is the same Chern--Simons vocabulary that governs the electron
mass formula of Paper~IV~\cite{Paper4}; as a mechanism for the oxygen
excess it is nonetheless excluded (Remark~\ref{P11:rem:oalpha_excess}).

The electron mass formula (Paper~IV, Theorem~8.4) %\ref{P4:thm:t_matrix_spoke}
is $m_e = \Lcond\cdot\vp^{20}\cdot e^{4\pi - T_{1/2}/Q}$
where $T_{1/2}/Q = 3/400$ is the WZW modular T-matrix correction
(independent of $\aem$, closing the $0.71\%$ spoke residual).
The remaining $0.22\%$ residual in that formula is the
one-loop QED $\overline{\mathrm{MS}}$-to-pole-mass conversion
$\aem/\pi = 0.2323\%$
(Paper~IV, Theorem~8.11); %\ref{P4:thm:msbar_pole}
incorporating it gives the fully corrected formula
$m_e = \Lcond\cdot\vp^{20}\cdot e^{4\pi-3/400-\aem/\pi}$
with a residual of $0.013\%$.

The oxygen IE correction $3\aem/(8\pi)$ and the electron mass
scheme conversion $\aem/\pi$ are related by
\begin{equation}
  \frac{3\aem}{8\pi} = \frac{3}{8} \cdot \frac{\aem}{\pi}:
  \label{P11:eq:kalpharelation}
\end{equation}
both are first-order in $\aem$ and both descend from the
WZW Chern--Simons phase $4\pi$ of Paper~IV via the
single-winding normalisation $2\pi$.
The candidate oxygen coefficient $3\aem/(8\pi)$ is smaller than the
electron-mass scheme term $\aem/\pi$ by the factor $3/8$; both are
$O(\aem)$ and share the Chern--Simons phase $4\pi$, though only the
electron-mass term is an established correction.
\end{remark}
\status{published}
\source{P11:rem:k8pi_origin}

\begin{remark}[Hund's rule and the H--O coincidence]
\label{P11:rem:hund_HO_coincidence}
The anomalous depression of oxygen's IE relative to nitrogen
($\IE(\mathrm{N}) = 14.534 > \IE(\mathrm{O}) = 13.618\,\mathrm{eV}$)
is a direct consequence of Hund's first rule: the half-filled $2p^3$
configuration of nitrogen is exchange-stabilised, giving an anomalously
high IE.  Adding a fourth 2p electron (oxygen) pairs it with an
existing $2p$ electron, incurring exchange penalty and
\emph{reducing} the IE below nitrogen's.
This reduction of exactly $0.916\,\mathrm{eV} = 0.0674\,\ERy$ brings
$\IE(\mathrm{O})$ to within $0.009\,\mathrm{eV}$ of $\ERy$
--- a coincidence of quantum many-body physics that is fixed once
$\aem$ and $m_e$ are set by the condensate.
\end{remark}
\status{published}
\source{P11:rem:hund_HO_coincidence}

\begin{remark}[Why water forms]
\label{P11:rem:why_water_forms}
The near-degeneracy $\IE_\mathrm{H} \approx \IE_\mathrm{O}$ means the
outermost electrons of H and O experience almost identical
binding energies.  This is a necessary (though not sufficient)
condition for covalent bond formation via orbital mixing:
the orbitals are at the same energy, so hybridisation
is energetically favourable.  The stability and geometry of
water is ultimately anchored to the Rydberg level being fixed
on the $\vp$-tower.
\end{remark}
\status{published}
\source{P11:rem:why_water_forms}

\begin{remark}[Alkali metal reactivity with water as a spiral-gap phenomenon]
\label{P11:rem:alkali_reactivity}
The $\vp$-spiral positions of the group-1 metals relative to the H--O
cluster quantify the thermodynamic driving force for the reaction
$2\,M + 2\,\mathrm{H_2O} \to 2\,\mathrm{MOH} + \mathrm{H_2}$.

All five stable alkali metals sit approximately one full spiral level
above the H--O cluster ($n \approx -12.10$):
\begin{center}
\begin{tabular}{lrrrl}
\toprule
Metal & IE (eV) & $n(\IE)$ & $|\text{gap}|$ & Reactivity with water \\
\midrule
Li & $5.392$ & $-11.140$ & $0.962$ & slow fizz, no flame \\
Na & $5.139$ & $-11.090$ & $1.012$ & vigorous, possible flame \\
K  & $4.341$ & $-10.915$ & $1.187$ & violent, bursts into flame \\
Rb & $4.177$ & $-10.875$ & $1.227$ & very violent, ignites \\
Cs & $3.894$ & $-10.802$ & $1.300$ & explosive \\
\bottomrule
\end{tabular}
\end{center}
The $|\text{gap}| \approx 1$ for all five elements reflects that the
alkali metal IE falls approximately one $\vp^2$-step below the
Rydberg level: $\IE_M \approx \ERy/\vp^2 \approx 5.2\,\mathrm{eV}$
for the group centre.
This common gap is what makes all alkali metals react with water ---
their outer electron is bound $\vp^2 \approx 2.6$ times more weakly
than the electrons in the H--O system, making donation thermodynamically
driven for every member of the group.

Within the group, the gap increases monotonically
Li $\to$ Na $\to$ K $\to$ Rb $\to$ Cs ($|\text{gap}|$ from
$0.962$ to $1.300$), and this ordering matches the observed reactivity
series exactly.
The framework prediction here is the same type as
Proposition~\ref{P11:prop:goldschmidt} for Goldschmidt pairs:
the spiral ordering of IEs predicts the chemical ordering of
thermodynamic driving forces, with no additional parameters.

The transition from vigorous (Na) to explosive (K, Rb, Cs) violence
is not explained by the spiral alone.
K, Rb, and Cs all have melting points below $100\,^{\circ}\mathrm{C}$
($63.4$, $39.3$, and $28.4\,^{\circ}\mathrm{C}$ respectively),
so they melt on contact with water and spread across the surface,
creating a kinetic amplification that is absent for Li
($180.5\,^{\circ}\mathrm{C}$) and Na ($97.8\,^{\circ}\mathrm{C}$).
The spiral quantifies the thermodynamic driving force; the melting
point governs the kinetic amplification that converts driving force
into violence.
\end{remark}
\status{published}
\source{P11:rem:alkali_reactivity}

%════════════════════════════════════════════════════════════════════
% Spiral Proximity and Goldschmidt — Paper XI
%════════════════════════════════════════════════════════════════════
\begin{theorem}[Spiral proximity predicts Goldschmidt compatibility]
\label{P11:thm:auc}
Let the binary label of an element pair $(A, B)$ be $1$ if the pair
is a known Goldschmidt isomorphous substitution pair~\cite{Goldschmidt1937,Burns1967}
and $0$ otherwise.
Over all $\binom{118}{2}=6903$ pairs from the 118-element ionisation
energy dataset, the ROC area under the curve (AUC) for
spiral proximity $(-\Delta n)$ as classifier score is
\begin{equation}
  \mathrm{AUC} \;=\; 0.857 \;>\; 0.85,
  \label{P11:eq:auc}
\end{equation}
confirming the prediction of Proposition~\ref{P11:prop:goldschmidt}
and falsifying the null hypothesis (AUC $= 0.5$) at high confidence.
At threshold $\Delta n < 0.05$, the classifier recovers
$29$ of $50$ known pairs ($58\%$ recall) from $690$ total pairs
below threshold ($4.2\%$ precision).
\end{theorem}
\status{published}
\source{P11:thm:auc}

\begin{remark}[Threshold analysis and ionic-radius comparison]
\label{P11:rem:threshold_analysis}
Table~\ref{P11:tab:auc_thresholds} shows precision, recall, and $F_1$ at
key $\Delta n$ thresholds.
The AUC of $0.857$ substantially exceeds the $\mathrm{AUC}\approx0.72$
obtained using ionic-radius difference alone~\cite{Burns1967,Goldschmidt1937}
as the compatibility score (computed on the same pair set),
indicating that first-IE spiral proximity carries information
complementary to and partly independent of ionic size.
The remaining $43\%$ recall gap (21 of 50 known pairs lie outside
$\Delta n < 0.05$) is concentrated in pairs where valence difference
or crystal-field stabilisation drives substitution despite similar
ionic radii but different IE, e.g.\ $(\mathrm{Ca}, \mathrm{Pb})$
($\Delta n = 0.200$) and $(\mathrm{Li}, \mathrm{Mg})$
($\Delta n = 0.363$).
\end{remark}
\status{published}
\source{P11:rem:threshold_analysis}

\begin{proposition}[Spiral proximity of Goldschmidt pairs]
\label{P11:prop:goldschmidt}
The five major Goldschmidt isomorphous substitution pairs
in mantle mineralogy --- (Mg, Fe), (Mg, Ni), (Fe, Co),
(Fe, Ni), (Ca, Sr) --- all satisfy $\Delta n < 0.05$ on the
$\vp$-spiral.
\end{proposition}
\status{published}
\source{P11:prop:goldschmidt}

\begin{remark}[Predictive power — confirmed]
\label{P11:rem:predictive_power_confirmed}
Proposition~\ref{P11:prop:goldschmidt} raises the falsifiable hypothesis that
spiral proximity ($\Delta n < \Delta_\mathrm{thr}$) should predict
Goldschmidt substitution compatibility better than chance.
Section~\ref{P11:sec:goldschmidt_auc} tests this hypothesis against the full
50-pair database~\cite{Goldschmidt1937,Burns1967} and finds
$\mathrm{AUC} = 0.857$ (Theorem~\ref{P11:thm:auc}), confirming the prediction.
\end{remark}
\status{published}
\source{P11:rem:predictive_power_confirmed}

\begin{conjecture}[$\vp$-spiral periodic table]
\label{P11:conj:spiral_pt}
The first ionisation energies of all 118 elements (90 naturally occurring
plus 28 synthetic, where IEs are experimentally measured or reliably estimated)
follow a distribution on the $\vp$-tower that is non-uniform:
they cluster at specific tower levels set by the
atomic-shell structure ($n=-10$ to $-13$), with the cluster
boundaries corresponding to noble gas IEs (which mark
``full-shell'' configurations with anomalously high IEs).
The clustering is ultimately anchored to the Rydberg level
$n = -12.102$ derived in Corollary~\ref{P11:cor:rydberg}.
\end{conjecture}
\status{open}
\source{P11:conj:spiral_pt}

%════════════════════════════════════════════════════════════════════
% Successive Ionisation Energies and the Multi-Layer Spiral — Paper XI
%════════════════════════════════════════════════════════════════════
\begin{proposition}[Master formula for successive IEs]
\label{P11:prop:master_formula}
In the hydrogenic approximation, the $k$-th ionisation energy of a
$Z$-electron atom satisfies
$\IE_k = Z_{\rm eff}(k)^2\,\ERy/n_{\rm qn}(k)^2$,
where $Z_{\rm eff}(k)$ is the effective nuclear charge seen by the
$k$-th electron removed and $n_{\rm qn}(k)$ its principal quantum
number.  Its tower level is
\begin{equation}
  \boxed{
    \nlvl{\IE_k}
    \;=\; \nlvl{\ERy}
          \;-\; \log_\vp Z_{\rm eff}(k)
          \;+\; \log_\vp n_{\rm qn}(k),
  }
  \label{P11:eq:master_formula}
\end{equation}
with $\nlvl{\ERy}=-12.102$ fixed by Corollary~\ref{P11:cor:rydberg}
and no new free parameters.
\end{proposition}
\status{published}
\source{P11:prop:master_formula}

\begin{remark}[The first-IE layer as the $k=1$ case]
\label{P11:rem:first_IE_special}
Setting $k=1$ recovers the placement of elements in
Section~\ref{P11:sec:elements}: each element sits at the tower level
set by its outermost $Z_{\rm eff}$ and $n_{\rm qn}$.
The master formula~\eqref{P11:eq:master_formula} extends this to
every electron in every atom, converting the single-layer spiral into
a multi-layer structure with one layer per successive IE.
\end{remark}
\status{published}
\source{P11:rem:first_IE_special}

\begin{proposition}[Shell-completion near-integer positions]
\label{P11:prop:shell_completion}
The following shell-completion ionisation energies land within
$|\Delta n|\leq 0.017$ of an integer tower level:
\begin{center}
\begin{tabular}{lrrrl}
\toprule
Transition & $\IE_k$ (eV) & $\nlvl{\IE_k}$ & $|\Delta n|$ & Shell \\
\midrule
Pb  IE$_3$   & $31.938$ & $-12.989$ & $0.011$ & 6$p$ stripped \\
Xe  IE$_3$   & $32.123$ & $-12.995$ & $0.005$ & 5$p$ third electron \\
Xe  IE$_1$   & $12.130$ & $-11.983$ & $0.017$ & noble-gas 5$p$ \\
\bottomrule
\end{tabular}
\end{center}
Ionisation energies are from NIST~\cite{NIST_ASD}.
\end{proposition}
\status{published}
\source{P11:prop:shell_completion}

\begin{remark}[Xenon as the integer anchor on successive layers]
\label{P11:rem:xenon_anchor}
The fact that both Xe IE$_1$ and Xe IE$_3$ land within $0.017$ of
consecutive integers ($-12$ and $-13$ respectively) is the multi-layer
generalisation of the noble-gas high-IE boundary noted in
Section~\ref{P11:sec:elements}: noble gases mark integer positions on
\emph{every} layer of the multi-layer spiral, not just the first.
\end{remark}
\status{published}
\source{P11:rem:xenon_anchor}

\begin{proposition}[Inner-shell near-degeneracy: O\,IE$_6$ and Na\,IE$_5$]
\label{P11:prop:inner_degen}
Oxygen's sixth ionisation energy and sodium's fifth satisfy
\begin{equation}
  \IE_6(\mathrm{O}) = 138.12\,\mathrm{eV},
  \qquad
  \IE_5(\mathrm{Na}) = 138.40\,\mathrm{eV},
  \label{P11:eq:inner_IEs}
\end{equation}
with tower levels $\nlvl{\IE_6(\mathrm{O})}=-14.5102$ and
$\nlvl{\IE_5(\mathrm{Na})}=-14.5123$, giving
\begin{equation}
  |\Delta n(\mathrm{O}_6,\,\mathrm{Na}_5)| = 0.0021,
  \label{P11:eq:inner_degen_val}
\end{equation}
comparable to the H--O outer-shell near-degeneracy ($|\Delta n|=0.0015$).
\end{proposition}
\status{published}
\source{P11:prop:inner_degen}

\begin{remark}[Spiral crossings and chemical equivalence at high oxidation state]
\label{P11:rem:spiral_crossing}
Oxygen and sodium are $\Delta n=1.02$ apart on the first-IE layer
($n(\mathrm{O,1})=-12.10$, $n(\mathrm{Na,1})=-11.09$).
At the fifth/sixth IE layer they converge to within $\Delta n=0.0021$.
This spiral crossing predicts that O$^{6+}$ and Na$^{5+}$
(both isoelectronic with He$^0$ after their $2s$ electron is exposed)
present the same binding energy to a further electron: they become
chemically equivalent at the extreme oxidation state.
More generally, Proposition~\ref{P11:prop:master_formula} predicts
crossings wherever two elements share $Z_{\rm eff}(k)/n_{\rm qn}(k)$
at some ionisation level $k$; these mark the oxidation states at
which the elements are Goldschmidt-substitutable in the sense of
Section~\ref{P11:sec:goldschmidt}.
\end{remark}
\status{published}
\source{P11:rem:spiral_crossing}

%════════════════════════════════════════════════════════════════════
% Labeled Equations — Paper XI
%════════════════════════════════════════════════════════════════════
\begin{equation}
  H_Z = \sum_{i=1}^Z \left(-\frac{\nabla_i^2}{2m_e} - \frac{Z\aem}{r_i}\right)
       + \sum_{i<j} \frac{\aem}{r_{ij}}.
  \label{P11:eq:HZ}
\end{equation}

\begin{equation}
  \frac{H_Z}{2\ERy}
  \;=\; \sum_{i=1}^Z
        \!\left(-\tfrac{1}{2}\nabla_i^2 - \frac{Z}{r_i}\right)
  \;+\; \sum_{i<j}\frac{1}{r_{ij}},
  \label{P11:eq:HZ_au}
\end{equation}

\begin{equation}
  K_{2p,2p} = \frac{3}{25}\,F^2(2p),
  \qquad
  F^2(2p)\big|_{Z_\mathrm{eff}=2} = \frac{17}{81}\,Z_\mathrm{eff}\,\ERy,
  \label{P11:eq:K2p2p_angular}
\end{equation}

\begin{equation}
  \frac{\IE(\mathrm{O})}{\ERy}
  \;=\; -2\,\frac{\varepsilon_{2p}}{2\ERy}
  \;\equiv\; R(\text{O}),
  \label{P11:eq:R_O}
\end{equation}

\begin{equation}
  Z_\mathrm{eff}(\mathrm{O})
  = n_\mathrm{qn}\sqrt{\IE_\mathrm{O}/\ERy}
  = 2\sqrt{13.618/13.606}
  = 2.001.
  \label{P11:eq:Zeff_O}
\end{equation}

\begin{equation}
  Z_\mathrm{eff}(\mathrm{O}) \;\equiv\; n_\mathrm{qn}\sqrt{\IE(\mathrm{O})/\ERy}
  \;=\; 2\sqrt{13.618/13.606} \;=\; 2.001.
  \label{P11:eq:Zeff_def}
\end{equation}

\begin{equation}
  -\varepsilon_{2p}
  \;=\; \frac{Z_\mathrm{eff}^2}{n_\mathrm{qn}^2}\,\ERy,
  \qquad
  Z_\mathrm{eff} \;\equiv\; Z - \sigma_\mathrm{tot},
  \label{P11:eq:eps_Zeff}
\end{equation}

\begin{equation}
  \vp^{2\Delta n}
  = \vp^{21.89}
  \approx 37\,558
  = \frac{2}{\aem^2}.
  \label{P11:eq:gap}
\end{equation}

\begin{equation}
  \Delta n(\mathrm{Fe,Cu}) = 0.0234, \quad
  \Delta n(\mathrm{Al,Cu}) = 0.265, \quad
  \Delta n(\mathrm{Al,Fe}) = 0.289.
\end{equation}

\begin{equation}
  n_{\mathrm{Al}} = -0.8532, \quad
  n_{\mathrm{Fe}} = -0.5646, \quad
  n_{\mathrm{Cu}} = -0.5880,
\end{equation}

\begin{equation}
  \boxed{\frac{n_{\mathrm{Al}}}{n_{\mathrm{Cu}}} \;=\; \frac{63}{24}
  \;=\; 2.625 \;\approx\; \vp^2 \;=\; 2.618\,,\quad
  0.27\%\ \text{error},}
\end{equation}

\begin{align}
  \nlvl{\ERy}
  &= -\frac{\ln(\aem^2 m_e / 2\,\Lcond)}{2\ln\vp}
   = -\frac{\ln(m_e/\Lcond)}{2\ln\vp}
     -\frac{\ln(\aem^2/2)}{2\ln\vp} \notag\\
  &= n(m_e) + \frac{\ln(2/\aem^2)}{2\ln\vp}.
  \label{P11:eq:nRy_proof}
\end{align}

\begin{equation}
  \sigma_\mathrm{Slater}
  = 2\times0.85\;(1s^2) + 5\times0.35\;(\text{other } n{=}2)
  = 3.45,
  \qquad
  Z_\mathrm{eff}^\mathrm{Slater} = 4.55 .
  \label{P11:eq:sigma_slater}
\end{equation}

\begin{equation}
  F^0(2p,2p) = \tfrac{93}{256}\,Z\,\ERy,
  \qquad
  F^2(2p,2p) = \tfrac{45}{256}\,Z\,\ERy,
  \label{P11:eq:slater_condon}
\end{equation}

\begin{equation}
  \IE_\mathrm{Slater}(\mathrm{O})
  \;=\; E(\mathrm{O}^+) - E(\mathrm{O})
  \;=\; 14.17\,\mathrm{eV}
  \;=\; 1.041\,\ERy,
  \qquad
  |\Delta n| = 0.041 .
  \label{P11:eq:IE_slater_total}
\end{equation}

\begin{equation}
  m_n \;=\; \Lcond\,\vp^{-2n}, \qquad n \in \mathbb{R}, \qquad
  \vp = \tfrac{1+\sqrt{5}}{2},
  \label{P11:eq:tower}
\end{equation}

\begin{equation}
  \boxed{
    Z_\mathrm{eff}(\mathrm{O}) = 2
    \;\text{ is a consequence of the condensate}
  },
  \label{P11:eq:Zeff_condensate_origin}
\end{equation}

